\BeleskeID{09_2-s020}
% element: 09_2-s020:e05
Daljim porastom ulaznog napona izlazni napon opada. P tranzistor ostaje u zasićenju, dok N tranzistor prelazi u omsku oblast. Ravnoteža struja pre zamene napona glasi
\[\frac{k_p}2(V_{GS,P}-V_{Tp})^2=\frac{k_n}2\left(2V_{DS,N}(V_{GS,N}-V_{Tn})-V_{DS,N}^2\right)\]
Za tačku gornjeg ulaznog praga diferenciranje i zamena nagiba daju
\[0=-k_n\left(2(V_I-V_{Tn})-V_O\right)+k_nV_O(2+1)\]
Druga jednačina je $V_I=V_{Tn}+2V_O$. Rezultat za gornji prag može se pisati i u obliku
\[V_{IH}=V_{Tn}+2(V_{DD}+V_{Tp})\sqrt{\frac{k_p}{3k_n}}=V_{Tn}+(V_{DD}+V_{Tp})\sqrt{\frac{4k_p}{3k_n}}\]
Oba zapisa su ista, jer se faktor dva unosi pod koren kao četiri.
